\documentclass{amsart}
\usepackage{amsmath,amssymb,amsthm,hyperref}
\newtheorem{theorem}{Theorem}
\newtheorem{lemma}{Lemma}
\newtheorem{proposition}{Proposition}
\theoremstyle{definition}
\newtheorem{definition}{Definition}
\theoremstyle{remark}
\newtheorem{remark}{Remark}
\begin{document}

\title{An explicit recursive description of the generalized Cantor set}
\maketitle

\section{Setup and notation}

Let $(\alpha_n)_{n\ge1}$ satisfy $0<\alpha_n<1$ for all $n$ and the admissibility
condition
\begin{equation}\label{eq:adm}
1-\sum_{n=1}^{\infty}2^{n-1}\alpha_n\ \ge\ 0 .
\end{equation}
The set $C=\bigcap_{n\ge0}C_n$ is built by the middle--deletion procedure of the
problem statement. Throughout, $\{0,1\}^n$ denotes the set of binary strings of
length $n$ (with $\{0,1\}^0=\{()\}$, the empty string), and $\{0,1\}^{\mathbb N}$
denotes the set of infinite binary sequences $\varepsilon=(\varepsilon_k)_{k\ge1}$.
For $\varepsilon\in\{0,1\}^{\mathbb N}$ and $n\ge0$, $\varepsilon|n:=(\varepsilon_1,\dots,\varepsilon_n)$
is its length-$n$ prefix.

We first record the side length of the stage--$n$ intervals.

\begin{lemma}\label{lem:len}
For every $n\ge0$, each of the $2^{n}$ closed intervals making up $C_n$ has the
same length
\begin{equation}\label{eq:Ln}
L_n\ :=\ \frac{1}{2^{n}}\Bigl(1-\sum_{k=1}^{n}2^{k-1}\alpha_k\Bigr),
\end{equation}
with the convention that the empty sum is $0$, so $L_0=1$. The lengths obey the
recursion
\begin{equation}\label{eq:rec}
L_n=\frac{L_{n-1}-\alpha_n}{2}\qquad(n\ge1).
\end{equation}
Moreover $L_n\ge0$ for all $n$, the sequence $(L_n)_{n\ge0}$ is nonincreasing, and
\begin{equation}\label{eq:Linf}
L_\infty\ :=\ \lim_{n\to\infty}L_n\ =\ 1-\sum_{k=1}^{\infty}2^{k-1}\alpha_k\ \ge\ 0 .
\end{equation}
\end{lemma}

\begin{proof}
We argue by induction. For $n=0$, $C_0=[0,1]$ is one interval of length $L_0=1$,
matching \eqref{eq:Ln}. Assume $C_{n-1}$ consists of $2^{n-1}$ closed intervals
each of length $L_{n-1}$. At stage $n$ each such interval has its open middle
subinterval of (absolute) length $\alpha_n$ removed; this leaves two closed
intervals, one on each side, of equal length $L_n=(L_{n-1}-\alpha_n)/2$, which is
\eqref{eq:rec}. Thus $C_n$ is a union of $2^{n}$ closed intervals of common length
$L_n$. Solving \eqref{eq:rec} with $L_0=1$: if \eqref{eq:Ln} holds for $n-1$ then
\[
L_n=\tfrac12 L_{n-1}-\tfrac12\alpha_n
=\tfrac{1}{2^{n}}\Bigl(1-\sum_{k=1}^{n-1}2^{k-1}\alpha_k\Bigr)-\tfrac{\alpha_n}{2}
=\tfrac{1}{2^{n}}\Bigl(1-\sum_{k=1}^{n}2^{k-1}\alpha_k\Bigr),
\]
proving \eqref{eq:Ln}. Since the procedure removes nondegenerate middles, each
$L_n$ is the length of an actual interval and hence $L_n\ge0$. By \eqref{eq:rec},
$L_{n-1}-L_n=\tfrac12(L_{n-1}+\alpha_n)\ge0$ because $L_{n-1}\ge0$ and
$\alpha_n>0$, so $(L_n)$ is nonincreasing; bounded below by $0$, it converges. Its
limit is $\lim_n 2^{-n}\bigl(1-\sum_{k=1}^n 2^{k-1}\alpha_k\bigr)$; the partial
sums $\sum_{k=1}^n 2^{k-1}\alpha_k$ increase to $\sum_{k=1}^\infty
2^{k-1}\alpha_k\le1$ by \eqref{eq:adm}, whence
$L_n\to 1-\sum_{k=1}^{\infty}2^{k-1}\alpha_k=L_\infty\ge0$, proving \eqref{eq:Linf}.
\end{proof}

It is convenient to introduce the \emph{step sizes}
\begin{equation}\label{eq:delta}
\delta_k\ :=\ L_{k-1}-L_k\ =\ \tfrac12\bigl(L_{k-1}+\alpha_k\bigr)\ >\ 0
\qquad(k\ge1),
\end{equation}
where the middle equality is \eqref{eq:rec} and positivity holds because
$L_{k-1}\ge0$ and $\alpha_k>0$. Telescoping \eqref{eq:delta} gives, for every
$n\ge0$,
\begin{equation}\label{eq:tail}
\sum_{k=n+1}^{\infty}\delta_k\ =\ L_n-L_\infty\ \le\ L_n,
\qquad\text{in particular}\qquad
\sum_{k=1}^{\infty}\delta_k=1-L_\infty<\infty .
\end{equation}

\section{Labelling the intervals; the address map}

We label the $2^n$ intervals of $C_n$ by binary strings of length $n$ compatibly
with the parent--child relation: when the interval labelled
$\varepsilon=(\varepsilon_1,\dots,\varepsilon_n)$ is split at stage $n+1$, its
\emph{left} child gets the label $\varepsilon0$ and its \emph{right} child gets
$\varepsilon1$. Write $I_\varepsilon=[x_\varepsilon,\,x_\varepsilon+L_n]$ for the
interval labelled $\varepsilon\in\{0,1\}^n$.

\begin{proposition}\label{prop:endpoint}
For every $n\ge0$ and every $\varepsilon=(\varepsilon_1,\dots,\varepsilon_n)\in\{0,1\}^n$,
\begin{equation}\label{eq:xeps}
x_\varepsilon\ =\ \sum_{k=1}^{n}\varepsilon_k\,\delta_k ,
\qquad\text{so}\qquad
I_\varepsilon=\Bigl[\sum_{k=1}^{n}\varepsilon_k\delta_k,\ \
\sum_{k=1}^{n}\varepsilon_k\delta_k+L_n\Bigr].
\end{equation}
Moreover the two children of $I_\varepsilon$ are
\begin{equation}\label{eq:children}
I_{\varepsilon0}=[\,x_\varepsilon,\ x_\varepsilon+L_{n+1}\,],
\qquad
I_{\varepsilon1}=[\,x_\varepsilon+\delta_{n+1},\ x_\varepsilon+\delta_{n+1}+L_{n+1}\,],
\end{equation}
and they are separated by the deleted open gap $(x_\varepsilon+L_{n+1},\,x_\varepsilon+\delta_{n+1})$ of length $\alpha_{n+1}>0$.
\end{proposition}

\begin{proof}
Induction on $n$. For $n=0$ the only string is empty and $x_{()}=0$, the left
endpoint of $C_0=[0,1]$. Assume \eqref{eq:xeps} for length $n$, and fix
$\varepsilon\in\{0,1\}^n$ with $I_\varepsilon=[x_\varepsilon,x_\varepsilon+L_{n}]$.
Deleting the open middle of length $\alpha_{n+1}$ leaves two closed intervals each
of length $L_{n+1}$: the left one occupies $[x_\varepsilon,x_\varepsilon+L_{n+1}]$,
then a gap of length $\alpha_{n+1}$, then the right one of length $L_{n+1}$; since
$L_{n+1}+\alpha_{n+1}+L_{n+1}=L_n$ by \eqref{eq:rec}, the right one is
$[x_\varepsilon+L_{n+1}+\alpha_{n+1},\,x_\varepsilon+L_n]$. Using
$\delta_{n+1}=L_n-L_{n+1}$ from \eqref{eq:delta} together with
$L_n=2L_{n+1}+\alpha_{n+1}$ from \eqref{eq:rec},
\[
L_{n+1}+\alpha_{n+1}=L_{n+1}+(L_n-2L_{n+1})=L_n-L_{n+1}=\delta_{n+1},
\]
which gives \eqref{eq:children} and the stated gap, of length
$\delta_{n+1}-L_{n+1}=\alpha_{n+1}>0$. Hence
\[
x_{\varepsilon0}=x_\varepsilon=\sum_{k=1}^{n}\varepsilon_k\delta_k,
\qquad
x_{\varepsilon1}=x_\varepsilon+\delta_{n+1}=\sum_{k=1}^{n+1}\varepsilon_k\delta_k
\ \ (\text{coordinate }n{+}1\text{ equal to }1),
\]
so $x_{\varepsilon'}=\sum_{k=1}^{n+1}\varepsilon'_k\delta_k$ for every length-$(n+1)$
string $\varepsilon'$, completing the induction.
\end{proof}

For $\varepsilon\in\{0,1\}^{\mathbb N}$ define the \emph{address value}
\begin{equation}\label{eq:phi}
\varphi(\varepsilon)\ :=\ \sum_{k=1}^{\infty}\varepsilon_k\,\delta_k ,
\end{equation}
which converges since $0\le\varepsilon_k\delta_k\le\delta_k$ and
$\sum_k\delta_k<\infty$ by \eqref{eq:tail}.

\section{Main theorem: the explicit description of $C$}

\begin{theorem}\label{thm:main}
Set $L_0=1$, $L_k=\tfrac12(L_{k-1}-\alpha_k)$, and
$\delta_k=L_{k-1}-L_k=\tfrac12(L_{k-1}+\alpha_k)$, so that explicitly
$L_n=2^{-n}\bigl(1-\sum_{k=1}^{n}2^{k-1}\alpha_k\bigr)$. Let $L_\infty=\lim_n L_n$
as in \eqref{eq:Linf}. Then the constructed set is
\begin{equation}\label{eq:general}
\boxed{\,C\ =\ \bigsqcup_{\varepsilon\in\{0,1\}^{\mathbb N}}
\bigl[\,\varphi(\varepsilon),\ \varphi(\varepsilon)+L_\infty\,\bigr],
\qquad
\varphi(\varepsilon)=\sum_{k=1}^{\infty}\varepsilon_k\,\delta_k,\,}
\end{equation}
a disjoint union of closed intervals of common length $L_\infty$ indexed by the
infinite binary sequences. In the (standard, measure-zero) case $L_\infty=0$ this
collapses to the explicit closed form
\begin{equation}\label{eq:main}
C\ =\ \Bigl\{\ \sum_{k=1}^{\infty}\varepsilon_k\,\delta_k\ :\
(\varepsilon_k)_{k\ge1}\in\{0,1\}^{\mathbb N}\ \Bigr\}
\ =\ \varphi\bigl(\{0,1\}^{\mathbb N}\bigr),
\end{equation}
and then $\varphi:\{0,1\}^{\mathbb N}\to C$ is a bijection.
\end{theorem}

\begin{proof}
Recall $C_n=\bigcup_{\varepsilon\in\{0,1\}^n}I_\varepsilon$ and
$C=\bigcap_{n\ge0}C_n$, with $I_\varepsilon$ as in \eqref{eq:xeps}.

\smallskip
\emph{Step 0 (nesting and disjointness).}
By \eqref{eq:children}, each child interval $I_{\varepsilon b}$ ($b\in\{0,1\}$) is
contained in its parent $I_\varepsilon$; hence for every
$\varepsilon\in\{0,1\}^{\mathbb N}$ the prefix intervals are nested,
$I_{\varepsilon|0}\supseteq I_{\varepsilon|1}\supseteq\cdots$. Also by
\eqref{eq:children} the two children of any interval are separated by a gap of
length $\alpha_{n+1}>0$, hence are disjoint; consequently any two distinct strings
$\varepsilon\ne\varepsilon'$ of the same length $n$ have $I_\varepsilon\cap
I_{\varepsilon'}=\varnothing$ (split at the first coordinate where they differ and
apply the gap separation to their common ancestor).

\smallskip
\emph{Step 1 (the prefix interval shrinks to a cell).}
Fix $\varepsilon\in\{0,1\}^{\mathbb N}$. The left endpoints
$x_{\varepsilon|n}=\sum_{k=1}^{n}\varepsilon_k\delta_k$ increase to
$\varphi(\varepsilon)$, and the right endpoints
$x_{\varepsilon|n}+L_n$ decrease: indeed passing from $I_{\varepsilon|n}$ to
$I_{\varepsilon|(n+1)}$ either keeps the left endpoint and right endpoint
$x_{\varepsilon|n}+L_{n+1}\le x_{\varepsilon|n}+L_n$ (case $\varepsilon_{n+1}=0$),
or moves the right endpoint to $x_{\varepsilon|n}+\delta_{n+1}+L_{n+1}
=x_{\varepsilon|n}+L_n$ (case $\varepsilon_{n+1}=1$, by \eqref{eq:children}); in
both cases the right endpoint does not increase. As the right endpoints equal
$x_{\varepsilon|n}+L_n\to\varphi(\varepsilon)+L_\infty$, the nested intersection is
the closed cell
\begin{equation}\label{eq:cell}
J_\varepsilon:=\bigcap_{n\ge0}I_{\varepsilon|n}
=\bigl[\varphi(\varepsilon),\ \varphi(\varepsilon)+L_\infty\bigr].
\end{equation}

\smallskip
\emph{Step 2 ($C=\bigcup_\varepsilon J_\varepsilon$).}
($\supseteq$) For any $\varepsilon$ and any $n$, $J_\varepsilon\subseteq
I_{\varepsilon|n}\subseteq C_n$; hence $J_\varepsilon\subseteq\bigcap_n C_n=C$.

($\subseteq$) Let $t\in C$. For each $n$, $t\in C_n$, so $t\in I_{\sigma^{(n)}}$
for a string $\sigma^{(n)}\in\{0,1\}^n$ that is \emph{unique} by the disjointness
in Step 0. If $\tau$ denotes the length-$n$ prefix of $\sigma^{(n+1)}$, then
$t\in I_{\sigma^{(n+1)}}\subseteq I_\tau$ (nesting), so uniqueness at level $n$
forces $\tau=\sigma^{(n)}$; thus the $\sigma^{(n)}$ are nested prefixes and define
a single $\varepsilon\in\{0,1\}^{\mathbb N}$ with $\varepsilon|n=\sigma^{(n)}$ for
all $n$. Hence $t\in\bigcap_n I_{\varepsilon|n}=J_\varepsilon$. Therefore
$C=\bigcup_{\varepsilon}J_\varepsilon$, which combined with \eqref{eq:cell} is
exactly the union in \eqref{eq:general}.

\smallskip
\emph{Step 3 (the cells are pairwise disjoint, giving $\sqcup$).}
Let $\varepsilon\ne\varepsilon'$ and let $m\ge1$ be the first index where they
differ; say $\varepsilon_m=0,\ \varepsilon'_m=1$ (the other case is symmetric).
Both prefix strings agree up to $m-1$, so $J_\varepsilon\subseteq I_{(\varepsilon|m)}
=I_{(\varepsilon|(m-1))0}$ and $J_{\varepsilon'}\subseteq I_{(\varepsilon'|m)}
=I_{(\varepsilon|(m-1))1}$, the two children of the common ancestor
$I_{\varepsilon|(m-1)}$. These children are disjoint by Step 0, hence
$J_\varepsilon\cap J_{\varepsilon'}=\varnothing$. This proves the union in
\eqref{eq:general} is disjoint, establishing \eqref{eq:general} in full
generality.

\smallskip
\emph{Step 4 (the case $L_\infty=0$).}
If $L_\infty=0$ then each cell $J_\varepsilon=\{\varphi(\varepsilon)\}$ is a single
point, so \eqref{eq:general} becomes
$C=\{\varphi(\varepsilon):\varepsilon\in\{0,1\}^{\mathbb N}\}
=\varphi(\{0,1\}^{\mathbb N})$, which is \eqref{eq:main}. Surjectivity of
$\varphi$ onto $C$ is then immediate. For injectivity, with $m$ the first
differing index ($\varepsilon_m=0,\varepsilon'_m=1$), \eqref{eq:tail} gives
\[
\varphi(\varepsilon')-\varphi(\varepsilon)
=\delta_m+\sum_{k=m+1}^{\infty}(\varepsilon'_k-\varepsilon_k)\delta_k
\ \ge\ \delta_m-\sum_{k=m+1}^{\infty}\delta_k
=\delta_m-(L_m-L_\infty).
\]
Using $\delta_m=L_{m-1}-L_m$ and $L_{m-1}=2L_m+\alpha_m$ (from
\eqref{eq:delta},\eqref{eq:rec}),
\[
\delta_m-(L_m-L_\infty)=(L_{m-1}-L_m)-L_m+L_\infty=L_{m-1}-2L_m+L_\infty
=\alpha_m+L_\infty\ \ge\ \alpha_m>0,
\]
so $\varphi(\varepsilon')>\varphi(\varepsilon)$; hence $\varphi$ is injective, thus
a bijection $\{0,1\}^{\mathbb N}\to C$.
\end{proof}

\begin{remark}
The estimate in Step 4 is valid for any $L_\infty\ge0$ and shows
$\varphi(\varepsilon')-\varphi(\varepsilon)\ge\alpha_m+L_\infty$ at the first
differing coordinate $m$: $\varphi$ is always strictly order preserving, and it
parametrizes the left endpoints of the cells in \eqref{eq:general}. When
$L_\infty>0$ (i.e. \eqref{eq:adm} is strict, $\sum_{k}2^{k-1}\alpha_k<1$), the
cells $[\varphi(\varepsilon),\varphi(\varepsilon)+L_\infty]$ are nondegenerate
intervals and $C$ has positive Lebesgue measure $L_\infty$; the genuine Cantor set
(measure $0$, totally disconnected, perfect) is the case $L_\infty=0$, i.e.
$\sum_{k=1}^{\infty}2^{k-1}\alpha_k=1$, where \eqref{eq:main} holds.
\end{remark}

\begin{remark}[Classical middle--thirds set]
Taking $\alpha_n=3^{-n}$ gives $\sum_{n\ge1}2^{n-1}\alpha_n=\tfrac13\sum_{n\ge1}(2/3)^{n-1}=1$,
so $L_\infty=0$, and one computes $\delta_k=L_{k-1}-L_k=2\cdot3^{-k}$. Then
\eqref{eq:main} reads
$C=\bigl\{\sum_{k\ge1}2\varepsilon_k 3^{-k}:\varepsilon_k\in\{0,1\}\bigr\}$, the
standard base-$3$ description of the Cantor middle--thirds set.
\end{remark}

\begin{remark}[Measure consistency]
From $|C_n|=2^nL_n=1-\sum_{k=1}^n2^{k-1}\alpha_k$ and $C=\bigcap_nC_n$ with
$C_{n}\subseteq C_{n-1}$, the Lebesgue measure of $C$ equals
$\lim_n|C_n|=1-\sum_{k=1}^{\infty}2^{k-1}\alpha_k=L_\infty$, consistent with
\eqref{eq:general} and the problem statement.
\end{remark}

\end{document}
